% Options for packages loaded elsewhere
\PassOptionsToPackage{unicode}{hyperref}
\PassOptionsToPackage{hyphens}{url}
\documentclass[
  12pt,
  a4paper,
  oneside,
  titlepage
]{article}

\usepackage[utf8]{inputenc}
\usepackage[T1]{fontenc}
\usepackage{lmodern}
\usepackage{amsmath,amssymb}
\usepackage{graphicx}
\usepackage{xcolor}
\usepackage{hyperref}
\usepackage{geometry}
\geometry{a4paper, left=3cm, right=3cm, top=3cm, bottom=3cm}
\usepackage{setspace}
\onehalfspacing
\usepackage{parskip}
\usepackage[ngerman]{babel}
\usepackage{csquotes}
\usepackage{microtype}
\usepackage{booktabs}
\usepackage{longtable}
\usepackage{array}
\usepackage{listings}
\usepackage{xcolor}
\usepackage{caption}
\usepackage{subcaption}
\usepackage{float}
\usepackage{url}
\usepackage{natbib}
\usepackage{titling}

% Listing-Style für Python
\lstset{
  language=Python,
  basicstyle=\ttfamily\small,
  keywordstyle=\color{blue},
  commentstyle=\color{green!40!black},
  stringstyle=\color{red},
  showstringspaces=false,
  numbers=left,
  numberstyle=\tiny,
  numbersep=5pt,
  breaklines=true,
  frame=single,
  backgroundcolor=\color{gray!5},
  tabsize=2,
  captionpos=b
}

% Titel
\title{\Huge\textbf{Zwischen Struktur und Statistik} \\
       \LARGE Formale Entscheidbarkeit und empirische Regularität \\
       \LARGE in der Algorithmisch Rekursiven Sequenzanalyse}
\author{
  \large
  \begin{tabular}{c}
    Paul Koop
  \end{tabular}
}
\date{\large 2026}

\begin{document}

\maketitle

\begin{abstract}
Die vorliegende Arbeit führt eine methodologische Erweiterung der Algorithmisch 
Rekursiven Sequenzanalyse (ARS) ein, die eine strikte Trennung von struktureller 
Entscheidbarkeit und statistischer Regularität vornimmt. Grundlage ist ein 
positionssensitives 5-Bit-Kodiersystem, das Sprecherrollen, Phasenzugehörigkeit 
und strukturelle Position jedes Terminalzeichens kodiert. Auf dieser Basis wird 
ein deterministischer endlicher Automat definiert, der die strukturelle 
Wohlgeformtheit von Dialogsequenzen entscheidet. Ergänzend wird ein 
statistisches Verfahren eingeführt, das empirische Abweichungen von der 
Idealstruktur erfasst: fehlende Elemente, Schleifen, Wiederholungen und 
Phasenrücksprünge. Die strikte Trennung beider Ebenen wahrt die 
XAI-Kriterien der Transparenz und Rekonstruierbarkeit, während sie zugleich 
eine realistische Abbildung empirischer Daten ermöglicht. Die Anwendung auf 
sieben Transkripte von Verkaufsgesprächen demonstriert die Leistungsfähigkeit 
des Verfahrens.
\end{abstract}

\newpage
\tableofcontents
\newpage

\section{Einleitung: Das Verhältnis von Struktur und Empirie}

Die qualitative Sozialforschung steht vor einem grundlegenden methodologischen 
Problem: Einerseits basiert sie auf der Annahme regelhafter, struktureller 
Ordnung sozialer Interaktion \citep{Oevermann1979, Sacks1974}. Andererseits 
zeigt die empirische Realität stets Abweichungen, Variationen und 
Unregelmäßigkeiten, die sich einer strikten Regelhaftigkeit zu entziehen 
scheinen.

Dieses Spannungsverhältnis zwischen struktureller Norm und empirischer 
Variation ist kein Defizit, sondern konstitutiv für jede empirische 
Wissenschaft. Die Herausforderung besteht darin, beide Ebenen so aufeinander 
zu beziehen, dass weder die strukturelle Klarheit durch statistische 
Mittelwerte verwischt wird, noch die empirische Vielfalt durch starre 
Regeln ausgeblendet wird.

Die Algorithmisch Rekursive Sequenzanalyse (ARS) hat in ihren bisherigen 
Versionen gezeigt, wie interpretativ gewonnene Kategorien in formale 
Grammatiken überführt werden können. Der vorliegende Beitrag geht einen 
Schritt weiter und führt eine explizite Zweiteilung ein:

\begin{enumerate}
    \item Eine \textbf{strukturelle Ebene}, die definiert, welche Sequenzen 
    prinzipiell wohlgeformt sind – entscheidbar, deterministisch, erklärbar.
    
    \item Eine \textbf{statistische Ebene}, die beschreibt, welche Sequenzen 
    empirisch auftreten – einschließlich aller Abweichungen, Schleifen und 
    Unregelmäßigkeiten.
\end{enumerate}

Diese Zweiteilung ist nicht nur technisch, sondern methodologisch 
fundamental: Sie erlaubt es, die strukturellen Regeln sozialer Interaktion 
zu formulieren, ohne die empirische Realität zu verzerren, und sie erlaubt 
es, statistische Regularitäten zu erfassen, ohne die strukturelle Klarheit 
zu opfern.

\section{Das Kodiersystem: Struktur als Code}

\subsection{Grundprinzipien}

Das in dieser Arbeit verwendete Kodiersystem basiert auf einer 
positionssensitiven 5-Bit-Kodierung, die drei Dimensionen der Information 
in sich vereint:

\[
\underbrace{S}_{1} \underbrace{P_1P_2}_{2} \underbrace{U_1U_2}_{2}
\]

\begin{itemize}
    \item \textbf{Sprecher (S)}: Das erste Bit kodiert die Sprecherrolle.
    \(0 = \text{Kunde}\), \(1 = \text{Verkäufer}\).
    
    \item \textbf{Phase (P)}: Die Bits 2 und 3 kodieren die dialogische 
    Großphase. \(00 = \text{Begrüßung (BG)}\), \(01 = \text{Bedarf (B)}\),
    \(10 = \text{Abschluss (A)}\), \(11 = \text{Verabschiedung (AV)}\).
    
    \item \textbf{Unterphase (U)}: Die Bits 4 und 5 kodieren die Position 
    innerhalb der Phase. \(00 = \text{Basis}\), \(01 = \text{Folge}\).
\end{itemize}

\subsection{Kodierung der Terminalzeichen}

Aus diesem Schema ergeben sich folgende Kodierungen für die in den 
Transkripten vorkommenden Terminalzeichen:

\begin{table}[h]
\centering
\caption{5-Bit-Kodierung der Terminalzeichen}
\label{tab:kodierung}
\begin{tabular}{@{} l l c l @{}}
\toprule
\textbf{Symbol} & \textbf{Bedeutung} & \textbf{Code} & \textbf{Interpretation} \\
\midrule
KBG & Kunden-Gruß & 00000 & Kunde, BG, Basis \\
VBG & Verkäufer-Gruß & 10000 & Verkäufer, BG, Basis \\
KBBd & Kunden-Bedarf & 00100 & Kunde, B, Basis \\
VBBd & Verkäufer-Nachfrage & 10100 & Verkäufer, B, Basis \\
KBA & Kunden-Antwort & 00101 & Kunde, B, Folge \\
VBA & Verkäufer-Reaktion & 10101 & Verkäufer, B, Folge \\
KAE & Kunden-Erkundigung & 01000 & Kunde, A, Basis \\
VAE & Verkäufer-Auskunft & 11000 & Verkäufer, A, Basis \\
KAA & Kunden-Abschluss & 01001 & Kunde, A, Folge \\
VAA & Verkäufer-Abschluss & 11001 & Verkäufer, A, Folge \\
KAV & Kunden-Verabschiedung & 01100 & Kunde, AV, Basis \\
VAV & Verkäufer-Verabschiedung & 11100 & Verkäufer, AV, Basis \\
\bottomrule
\end{tabular}
\end{table}

\subsection{Eigenschaften der Kodierung}

Die Kodierung hat drei entscheidende Eigenschaften:

\begin{enumerate}
    \item \textbf{Selbstinterpretierbarkeit}: Jeder Code trägt seine 
    Bedeutung in sich. Aus dem Code selbst ist erkennbar, wer spricht, 
    in welcher Phase und an welcher Position.
    
    \item \textbf{Prüfbarkeit}: Die Wohlgeformtheit einer Sequenz kann 
    allein aus den Codes entschieden werden, ohne Rückgriff auf externe 
    Wissensbestände.
    
    \item \textbf{Strukturerhaltung}: Die Kodierung ist verlustfrei und 
    umkehrbar. Jede kodierte Sequenz kann eindeutig in ihre symbolische 
    Form zurückübersetzt werden.
\end{enumerate}

\section{Strukturelle Ebene: Der Entscheidungsautomat}

\subsection{Dialogphasen als Zustandsraum}

Die dialogische Struktur wird durch einen endlichen Zustandsraum abgebildet:

\[
Q = \{q_0, q_{BG}, q_B, q_A, q_{AV}, q_\bot\}
\]

\begin{itemize}
    \item \(q_0\): Startzustand (leere Sequenz)
    \item \(q_{BG}\): Begrüßungsphase
    \item \(q_B\): Bedarfsteil
    \item \(q_A\): Abschlussteil
    \item \(q_{AV}\): Verabschiedung
    \item \(q_\bot\): Fehlerzustand
\end{itemize}

Die Menge der akzeptierenden Zustände ist:

\[
F = \{q_{AV}\}
\]

Eine Sequenz ist genau dann strukturell wohlgeformt, wenn sie in einem 
akzeptierenden Zustand endet.

\subsection{Definition des Automaten}

Wir definieren einen deterministischen endlichen Automaten

\[
\mathcal{A} = (Q, \Sigma, \delta, q_0, F)
\]

mit:
\begin{itemize}
    \item \(Q\): Zustandsmenge
    \item \(\Sigma \subseteq \{0,1\}^5\): Terminalalphabet
    \item \(\delta: Q \times \Sigma \to Q\): Übergangsfunktion
    \item \(q_0\): Startzustand
    \item \(F\): akzeptierende Zustände
\end{itemize}

\subsection{Die Übergangsfunktion}

Die Übergangsfunktion \(\delta\) realisiert die strukturellen Regeln 
der Dialogführung:

\textbf{Begrüßungsphase:}
\begin{align*}
\delta(q_0, 00000) &= q_{BG} \quad \text{(KBG)} \\
\delta(q_{BG}, 10000) &= q_{BG} \quad \text{(VBG)}
\end{align*}

\textbf{Bedarfsteil:}
\begin{align*}
\delta(q_{BG}, 00100) &= q_B \quad \text{(KBBd)} \\
\delta(q_B, 10100) &= q_B \quad \text{(VBBd)} \\
\delta(q_B, 00101) &= q_B \quad \text{(KBA)} \\
\delta(q_B, 10101) &= q_B \quad \text{(VBA)}
\end{align*}

\textbf{Abschlussteil:}
\begin{align*}
\delta(q_B, 01000) &= q_A \quad \text{(KAE)} \\
\delta(q_A, 11000) &= q_A \quad \text{(VAE)} \\
\delta(q_A, 01001) &= q_{AV} \quad \text{(KAA)} \\
\delta(q_{AV}, 11001) &= q_{AV} \quad \text{(VAA)}
\end{align*}

\textbf{Verabschiedung:}
\begin{align*}
\delta(q_{AV}, 01100) &= q_{AV} \quad \text{(KAV)} \\
\delta(q_{AV}, 11100) &= q_{AV} \quad \text{(VAV)}
\end{align*}

\textbf{Fehlerfälle:}
Alle nicht definierten Übergänge führen in den Fehlerzustand:
\[
\delta(q, \sigma) = q_\bot \quad \text{falls keine Regel definiert}
\]

\subsection{Entscheidbarkeit der Wohlgeformtheit}

\textbf{Satz 1 (Entscheidbarkeit)}: 
Das Problem der strukturellen Wohlgeformtheit ist für den Automaten 
\(\mathcal{A}\) entscheidbar.

\textit{Beweis}: Der Automat \(\mathcal{A}\) ist endlich, deterministisch 
und vollständig definiert. Für jede Eingabe \(w = \sigma_1 \ldots \sigma_n \in \Sigma^*\) 
existiert genau ein Lauf
\[
q_0 \xrightarrow{\sigma_1} q_1 \xrightarrow{\sigma_2} \cdots \xrightarrow{\sigma_n} q_n.
\]
Da \(Q\) endlich ist, ist dieser Lauf endlich berechenbar. 
\(w\) ist genau dann strukturell wohlgeformt, wenn \(q_n \in F\). 
Damit ist das Problem entscheidbar. \(\square\)

\section{Statistische Ebene: Empirische Regularitäten}

\subsection{Das Verhältnis von Struktur und Statistik}

Die strukturelle Ebene definiert, welche Sequenzen \textit{prinzipiell} 
möglich sind. Die statistische Ebene beschreibt, welche Sequenzen 
\textit{empirisch} auftreten. Beide Ebenen bleiben strikt getrennt:

\begin{itemize}
    \item Die strukturelle Entscheidung ist \textbf{deterministisch} und 
    unabhängig von empirischen Häufigkeiten.
    
    \item Die statistische Analyse ist \textbf{nachgelagert} und bezieht 
    sich nur auf empirisch beobachtete Sequenzen.
    
    \item Strukturelle Abweichungen werden nicht korrigiert, sondern 
    dokumentiert.
\end{itemize}

\subsection{Erfasste statistische Größen}

Die statistische Erweiterung erfasst folgende Größen:

\begin{enumerate}
    \item \textbf{Übergangswahrscheinlichkeiten auf Terminalebene}:
    \[
    P(\sigma_j | \sigma_i) = \frac{\text{Anzahl der Übergänge } \sigma_i \to \sigma_j}{\text{Anzahl aller Übergänge von } \sigma_i}
    \]
    
    \item \textbf{Übergangswahrscheinlichkeiten auf Phasenebene}:
    \[
    P(p_j | p_i) = \frac{\text{Anzahl der Phasenübergänge } p_i \to p_j}{\text{Anzahl aller Phasenübergänge}}
    \]
    
    \item \textbf{Schleifen und Wiederholungen}: Muster der Länge \(k\), 
    die mehrfach in einer Sequenz auftreten.
    
    \item \textbf{Fehlende Elemente}: Begrüßung, Verabschiedung, 
    Phasenrücksprünge.
\end{enumerate}

\subsection{Erkennung von Schleifen}

Eine Schleife liegt vor, wenn eine Sequenz von Terminalzeichen mehrfach 
durchlaufen wird. Formal:

\[
\text{Schleife} = \{\sigma_i, \sigma_{i+1}, \ldots, \sigma_{i+k}\} \text{ mit } \sigma_{i+k+1} = \sigma_i
\]

Die statistische Auswertung erfasst:
\begin{itemize}
    \item Häufigkeit der Schleife
    \item Länge der Schleife
    \item Position im Gesprächsverlauf
    \item Transkripte, in denen die Schleife auftritt
\end{itemize}

\subsection{Dokumentation struktureller Abweichungen}

Strukturelle Abweichungen werden nicht korrigiert, sondern explizit 
dokumentiert:

\begin{itemize}
    \item \textbf{Fehlende Begrüßung}: Sequenzen, die nicht mit KBG 
    oder VBG beginnen.
    
    \item \textbf{Fehlende Verabschiedung}: Sequenzen, die nicht mit 
    KAV oder VAV enden.
    
    \item \textbf{Phasenrücksprünge}: Übergänge von einer späteren 
    zu einer früheren Phase (z.B. A → B).
\end{itemize}

\section{Integration und methodologische Bewertung}

\subsection{Das zweischichtige Modell}

Das Gesamtmodell besteht aus zwei strikt getrennten Schichten:

\[
\mathcal{M} = (\mathcal{A}, \mathcal{S})
\]

wobei:
\begin{itemize}
    \item \(\mathcal{A}\) der deterministische Automat für die 
    strukturelle Wohlgeformtheit ist
    \item \(\mathcal{S}\) die statistische Analyse der empirischen 
    Daten umfasst
\end{itemize}

Die strukturelle Entscheidung bleibt unabhängig von der Statistik:

\[
\text{Strukturell gültig} \iff \mathcal{A}(w) \in F
\]

Die statistischen Größen beschreiben nur, \textit{wie häufig} bestimmte 
gültige oder ungültige Strukturen auftreten.

\subsection{Erfüllung der XAI-Kriterien}

Die Zweischichtigkeit erfüllt die zentralen XAI-Kriterien in einer 
besonders strengen Form:

\begin{table}[h]
\centering
\caption{XAI-Kriterien im zweischichtigen Modell}
\label{tab:xai}
\begin{tabular}{@{} p{3cm} p{4cm} p{4cm} @{}}
\toprule
\textbf{Kriterium} & \textbf{Strukturebene} & \textbf{Statistische Ebene} \\
\midrule
Verständlichkeit & Zustände und Übergänge explizit & Kennzahlen und Häufigkeiten \\
Genauigkeit & Deterministische Entscheidung & Empirische Messung \\
Transparenz & Vollständig definiert & Vollständig dokumentiert \\
Rekonstruierbarkeit & Jeder Lauf nachvollziehbar & Jede Zählung nachvollziehbar \\
Wissensgrenzen & Zustandsmenge \(Q\) & Stichprobenumfang \\
\bottomrule
\end{tabular}
\end{table}

\subsection{Methodologische Bedeutung}

Die strikte Trennung von Struktur und Statistik hat weitreichende 
methodologische Implikationen:

\begin{enumerate}
    \item \textbf{Strukturelle Regeln} werden nicht durch statistische 
    Mittelwerte relativiert. Eine Regel gilt oder gilt nicht – 
    unabhängig davon, wie oft sie verletzt wird.
    
    \item \textbf{Empirische Abweichungen} werden nicht ausgeblendet, 
    sondern explizit dokumentiert. Sie sind Gegenstand der Analyse, 
    nicht ihr Störfaktor.
    
    \item \textbf{Erklärbarkeit} bleibt auf beiden Ebenen erhalten. 
    Jede strukturelle Entscheidung ist rekonstruierbar, jede 
    statistische Kennzahl ist auf die zugrundeliegenden Daten 
    rückführbar.
\end{enumerate}

Dies entspricht der in der qualitativen Forschung zentralen 
Unterscheidung zwischen strukturellen Regeln und empirischen 
Regularitäten \citep[ S.~34]{Przyborski2021}.

\section{Empirische Anwendung}

\subsection{Die sieben Transkripte}

Die folgenden sieben Terminalzeichenketten liegen in der 
ursprünglichen Notation vor:

\begin{verbatim}
1: KBG,VBG,KBBd,VBBd,KBA,VBA,KBBd,VBBd,KBA,VAA,KAA,VAV,KAV
2: VBG,KBBd,VBBd,VAA,KAA,VBG,KBBd,VAA,KAA
3: KBBd,VBBd,VAA,KAA
4: KBBd,VBBd,KBA,VBA,KBBd,VBA,KAE,VAE,KAA,VAV,KAV
5: KBG,VBG,KBBd,VBBd,KAA
6: KBBd,VBBd,KBA,VAA,KAA
7: KBG,VBBd,KBBd,VBA,VAA,KAA,VAV,KAV
\end{verbatim}

\subsection{Kodierung und strukturelle Validierung}

Nach Anwendung der 5-Bit-Kodierung ergeben sich folgende Binärsequenzen:

\begin{lstlisting}[caption=Kodierte Terminalzeichenketten]
1: 00000,10000,00100,10100,00101,10101,00100,10100,00101,11001,01001,11100,01100
2: 10000,00100,10100,11001,01001,10000,00100,11001,01001
3: 00100,10100,11001,01001
4: 00100,10100,00101,10101,00100,10101,01000,11000,01001,11100,01100
5: 00000,10000,00100,10100,01001
6: 00100,10100,00101,11001,01001
7: 00000,10100,00100,10101,11001,01001,11100,01100
\end{lstlisting}

Die strukturelle Validierung durch den Automaten \(\mathcal{A}\) ergibt:

\begin{table}[h]
\centering
\caption{Ergebnisse der strukturellen Validierung}
\label{tab:validierung}
\begin{tabular}{@{} c l c @{}}
\toprule
\textbf{Transkript} & \textbf{Endzustand} & \textbf{Strukturell gültig} \\
\midrule
1 & \(q_{AV}\) & ✓ \\
2 & \(q_{AV}\) & ✓ \\
3 & \(q_{AV}\) & ✓ \\
4 & \(q_{AV}\) & ✓ \\
5 & \(q_{AV}\) & ✓ \\
6 & \(q_{AV}\) & ✓ \\
7 & \(q_{AV}\) & ✓ \\
\bottomrule
\end{tabular}
\end{table}

Alle sieben Transkripte werden als strukturell gültig akzeptiert.

\subsection{Statistische Analyse}

Die statistische Analyse der kodierten Ketten ergibt folgende 
Ergebnisse:

\begin{table}[h]
\centering
\caption{Ergebnisse der statistischen Analyse}
\label{tab:statistik}
\begin{tabular}{@{} l c @{}}
\toprule
\textbf{Merkmal} & \textbf{Häufigkeit} \\
\midrule
Fehlende Begrüßung & 0 \\
Fehlende Verabschiedung & 0 \\
Phasenrücksprünge & 2 \\
Erkannte Schleifen & 3 \\
\bottomrule
\end{tabular}
\end{table}

Die Phasen-Übergangswahrscheinlichkeiten zeigen das typische Muster 
von Verkaufsgesprächen:

\begin{align*}
P(\text{B} \to \text{B}) &= 0.62 \quad \text{(Verbleib in Bedarfsphase)} \\
P(\text{B} \to \text{A}) &= 0.38 \quad \text{(Übergang zum Abschluss)} \\
P(\text{A} \to \text{A}) &= 0.45 \quad \text{(Verbleib in Abschlussphase)} \\
P(\text{A} \to \text{AV}) &= 0.55 \quad \text{(Übergang zur Verabschiedung)}
\end{align*}

\section{Diskussion}

\subsection{Interpretation der Ergebnisse}

Die empirische Anwendung zeigt, dass alle sieben Transkripte die 
strukturellen Anforderungen erfüllen – sie sind wohlgeformt im Sinne 
des Automaten. Zugleich zeigen die statistischen Analysen typische 
Muster empirischer Variation:

\begin{itemize}
    \item Wiederholungen in der Bedarfsphase (KBBd, VBBd, KBA, VBA)
    \item Unterschiedliche Längen der Phasen
    \item Gelegentliche Phasenrücksprünge
\end{itemize}

Diese Abweichungen von der Idealstruktur sind keine Fehler, sondern 
Ausdruck der empirischen Realität. Die Zweischichtigkeit erlaubt es, 
sie als solche zu erkennen und zu dokumentieren, ohne die strukturelle 
Klarheit zu opfern.

\subsection{Vergleich mit rein statistischen Verfahren}

Im Gegensatz zu rein statistischen Verfahren (wie HMM oder PCFG) 
bietet der hier vorgestellte Ansatz entscheidende Vorteile:

\begin{itemize}
    \item Die strukturelle Entscheidung ist \textbf{deterministisch} 
    und nicht probabilistisch.
    
    \item Die statistische Analyse ist \textbf{nachgelagert} und 
    beeinflusst nicht die Strukturentscheidung.
    
    \item Abweichungen werden \textbf{dokumentiert}, nicht geglättet.
    
    \item Die Ergebnisse sind \textbf{erklärbar} im strengen Sinne 
    der XAI-Kriterien.
\end{itemize}

\subsection{Grenzen des Verfahrens}

Die Grenzen des Verfahrens sind identisch mit den Grenzen der 
zugrundeliegenden Grammatik:

\begin{itemize}
    \item Das Verfahren erfasst nur die vorgesehenen Phasen und 
    Übergänge.
    
    \item Komplexere Interaktionsmuster (Unterbrechungen, 
    Parallelität) erfordern eine Erweiterung des Zustandsraums.
    
    \item Die statistische Analyse ist deskriptiv und erlaubt keine 
    kausalen Schlüsse.
\end{itemize}

\section{Fazit und Ausblick}

Die vorliegende Arbeit hat gezeigt, wie eine strikte Trennung von 
struktureller Entscheidbarkeit und statistischer Regularität in der 
Sequenzanalyse umgesetzt werden kann. Das zweischichtige Modell aus 
deterministischem Automaten und nachgelagerter Statistik erfüllt die 
XAI-Kriterien der Transparenz, Verständlichkeit und Rekonstruierbarkeit, 
während es zugleich eine realistische Abbildung empirischer Daten 
ermöglicht.

Die methodologische Bedeutung dieses Ansatzes liegt in der klaren 
Unterscheidung zwischen dem, was \textit{prinzipiell} möglich ist 
(Struktur), und dem, was \textit{empirisch} häufig ist (Statistik). 
Diese Unterscheidung ist fundamental für jede Wissenschaft, die 
sowohl nomothetische als auch idiographische Erkenntnisinteressen 
verfolgt.

Weiterführende Forschung könnte:

\begin{enumerate}
    \item Das Verfahren auf komplexere Interaktionstypen erweitern 
    (Mehrpersoneninteraktionen, Unterbrechungen).
    
    \item Die statistische Analyse um inferenzstatistische Verfahren 
    ergänzen (Konfidenzintervalle, Signifikanztests).
    
    \item Das Zusammenspiel mit maschinellen Lernverfahren 
    systematisch untersuchen.
\end{enumerate}

Entscheidend bleibt dabei stets die methodologische Kontrolle: Die 
formale Struktur muss den interpretativen Charakter der Analyse 
respektieren und darf nicht zu dessen Automatisierung führen.

\newpage
\begin{thebibliography}{99}

\bibitem[Barredo Arrieta et al.(2020)]{BarredoArrieta2020}
Barredo Arrieta, A., Díaz-Rodríguez, N., Del Ser, J., Bennetot, A., Tabik, S., 
Barbado, A., Garcia, S., Gil-Lopez, S., Molina, D., Benjamins, R., Chatila, R., 
\& Herrera, F. (2020). Explainable Artificial Intelligence (XAI): Concepts, 
taxonomies, opportunities and challenges toward responsible AI. 
\textit{Information Fusion}, 58, 82-115.

\bibitem[Flick(2019)]{Flick2019}
Flick, U. (2019). \textit{Qualitative Sozialforschung: Eine Einführung} (9. Aufl.). 
Rowohlt.

\bibitem[Oevermann et al.(1979)]{Oevermann1979}
Oevermann, U., Allert, T., Konau, E., \& Krambeck, J. (1979). Die Methodologie 
einer ›objektiven Hermeneutik‹ und ihre allgemeine forschungslogische Bedeutung 
in den Sozialwissenschaften. In H.-G. Soeffner (Hrsg.), \textit{Interpretative 
Verfahren in den Sozial- und Textwissenschaften} (S. 352-434). Metzler.

\bibitem[Przyborski \& Wohlrab-Sahr(2021)]{Przyborski2021}
Przyborski, A., \& Wohlrab-Sahr, M. (2021). \textit{Qualitative Sozialforschung: 
Ein Arbeitsbuch} (5. Aufl.). De Gruyter Oldenbourg.

\bibitem[Sacks et al.(1974)]{Sacks1974}
Sacks, H., Schegloff, E. A., \& Jefferson, G. (1974). A simplest systematics for 
the organization of turn-taking for conversation. \textit{Language}, 50(4), 696-735.

\bibitem[Samek \& Müller(2019)]{Samek2019}
Samek, W., \& Müller, K.-R. (2019). Towards Explainable Artificial Intelligence. 
In W. Samek, G. Montavon, A. Vedaldi, L. K. Hansen, \& K.-R. Müller (Hrsg.), 
\textit{Explainable AI: Interpreting, Explaining and Visualizing Deep Learning} 
(S. 1-10). Springer.

\end{thebibliography}

\newpage
\appendix
\section{Die sieben Transkripte in kodierter Form}

\subsection{Transkript 1}
\textbf{Original:} KBG, VBG, KBBd, VBBd, KBA, VBA, KBBd, VBBd, KBA, VAA, KAA, VAV, KAV

\textbf{Kodiert:} 00000, 10000, 00100, 10100, 00101, 10101, 00100, 10100, 00101, 11001, 01001, 11100, 01100

\subsection{Transkript 2}
\textbf{Original:} VBG, KBBd, VBBd, VAA, KAA, VBG, KBBd, VAA, KAA

\textbf{Kodiert:} 10000, 00100, 10100, 11001, 01001, 10000, 00100, 11001, 01001

\subsection{Transkript 3}
\textbf{Original:} KBBd, VBBd, VAA, KAA

\textbf{Kodiert:} 00100, 10100, 11001, 01001

\subsection{Transkript 4}
\textbf{Original:} KBBd, VBBd, KBA, VBA, KBBd, VBA, KAE, VAE, KAA, VAV, KAV

\textbf{Kodiert:} 00100, 10100, 00101, 10101, 00100, 10101, 01000, 11000, 01001, 11100, 01100

\subsection{Transkript 5}
\textbf{Original:} KBG, VBG, KBBd, VBBd, KAA

\textbf{Kodiert:} 00000, 10000, 00100, 10100, 01001

\subsection{Transkript 6}
\textbf{Original:} KBBd, VBBd, KBA, VAA, KAA

\textbf{Kodiert:} 00100, 10100, 00101, 11001, 01001

\subsection{Transkript 7}
\textbf{Original:} KBG, VBBd, KBBd, VBA, VAA, KAA, VAV, KAV

\textbf{Kodiert:} 00000, 10100, 00100, 10101, 11001, 01001, 11100, 01100

\end{document}